\section*{Cluster A Step 37: Uniqueness Stress-Test}

\paragraph{Grade.}
Finite-carrier stress-test. The computation tests whether the Step-36 uniqueness claim depends on a shape-flavored factor-local rule. It does not derive content or values.

\paragraph{Deflationary Correction.}
The old factor-local rule required an untouched active factor and therefore rejected single-factor structures by construction. Step 37 drops that rule and computes unbroken non-abelian subgroups after scalar breaking, including partial breaking inside a factor.

\paragraph{Precise Test.}
The Step-35 residual is rederived as \(12\). The precise unbroken-subgroup test leaves all \(12\): \(8\) in \(\texttt{2|3}\) and \(4\) in \(\texttt{4}\). Thus the single-factor family survives the precise neutral test.

\paragraph{Neutral Currency Relation.}
The scalar-stage currency tuple is the scalar witness component dimension, active-route count, and charge magnitude. The minimum currency selects \(8\) survivors, all in \(\texttt{2|3}\). The Step-14 reference support passes but is not unique.

\paragraph{Verdict.}
\(\texttt{GENUINE\_NEUTRAL\_STRUCTURE\_UNIQUENESS\_WITH\_CONTENT\_LIMIT}\). Structure-level uniqueness survives only after replacing the factor-local rule with a neutral scalar-stage currency relation. Content uniqueness remains unresolved.

\paragraph{Boundary.}
This is a finite structural stress-test, not a physical derivation. The remaining residual is a content boundary.
